\section{The coefficient, kernel and weight-point constructions for preuniformisation}
\label{sec:faithful-preuniformisation}
This section follows the original complete-local preuniformisation
construction.  It proves the coefficient uniqueness and gluing,
the finite-kernel local argument and the actual passage from compatible
weight points to the datum \(\mathrm{PreU}\).
The existence and algebraisation of the formal quotient tower are
separate producers, specified at the end.  None of the lemmas below
asserts that an arbitrary formal torus or an arbitrary rational
trivialisation is the tower of the given family.

\subsection{The normalized even cubical unit}
Let \(S\) be nonempty and connected, and let \(G/S\) be a nonempty
connected commutative group
with the actual torus quotient
\[
       1\longrightarrow T_X\xrightarrow{i}G\xrightarrow{q}B
                      \longrightarrow1.
 \tag{PUF1}\label{eq:preunif-torus-extension}
\]
Here \(X\) is a free abelian group, \(T_X\) is its split torus,
\(B/S\) is an abelian scheme, and \(q\) is faithfully flat, surjective
and affine, with its stated kernel square.
Retain the usual abelian-scheme function theorem in its full
base-change form: for every \(Z/S\), the pullback map
\(\Gamma(Z,\mathcal O)\to\Gamma(B\times_S Z,\mathcal O)\)
is an isomorphism. Thus arbitrary global functions on that
pullback, not only a preselected fibre-constant class, come from \(Z\).
All point identities in this subsection are functorial in the test
scheme over \(S\), including nonreduced test schemes.

\begin{srlemma}[A character family is constant on a connected parameter space]
\label{lem:preunif-character-family}
Let \(Y/S\) be nonempty and connected and let
\(\rho\in\Gamma(T_X\times_S Y,\mathcal O)^\times\) be multiplicative
in the torus variable.  There is a unique \(\chi\in X\) for which
\(\rho(t,y)=\chi(t)\) on every test scheme.
\end{srlemma}
\begin{proof}
Locally on an affine open of \(Y\), write the group-like function in
the group algebra as a finite sum \(\sum_\chi a_\chi e^\chi\).
The equality between comultiplication of this sum and its tensor
square says
\[
       a_\chi^2=a_\chi,\qquad
       a_\chi a_{\chi'}=0\ (\chi\ne\chi'),\qquad
       \sum_\chi a_\chi=1.
\]
Thus its coefficients are a partition of the parameter space by
open-and-closed sets, on each of which a single coefficient is one.
This argument uses idempotents, so it does not presume the base is
reduced.  Uniqueness of the group-algebra coefficients identifies the
characters on intersecting opens.  The resulting locally constant
map \(Y\to X\) is constant because \(Y\) is connected.
This is the coefficient content of
\srcid{existsUnique_character_of_mul}.
\end{proof}

\begin{srlemma}[The actual normalized even cubical unit is one]
\label{lem:preunif-even-cubic-unit}
Let \(u\in\Gamma(G,\mathcal O)^\times\).  Suppose
\[
 \begin{gathered}
 u(1)=1,\qquad u(x^{-1})=u(x),\\
 u(xyz)u(x)u(y)u(z)=u(xy)u(xz)u(yz)
 \quad\hbox{for every test scheme and }x,y,z.
 \end{gathered}
 \tag{PUF2}\label{eq:preunif-even-cube}
\]
Then \(u=1\).
\end{srlemma}
\begin{proof}
Define on \(T_X\times_S G\)
\[
       \rho(t,g)=
            \frac{u(i(t)g)}{u(i(t))u(g)}.
\]
Apply the cubical identity to \(i(t),i(t'),g\).
Cancellation gives
\(\rho(tt',g)=\rho(t,g)\rho(t',g)\).
Lemma~\ref{lem:preunif-character-family} gives one character \(\chi\)
independent of \(g\).  At \(g=1\) the ratio is one, so \(\chi=0\).
Consequently
\[
                       u(i(t)g)=u(i(t))u(g).
 \tag{PUF3}\label{eq:preunif-split-unit}
\]
Restriction of this equality to \(g=i(t')\) makes \(u\circ i\) a
character \(\psi\) of \(T_X\).  Evenness gives \(\psi=-\psi\).
The character lattice is torsion free, so \(2\psi=0\) implies
\(\psi=0\).  Equation~\eqref{eq:preunif-split-unit} therefore makes
\(u\) invariant under the actual torus kernel action.

Faithfully flat descent now gives a function \(a\) on \(B\):
on \(G\times_B G\), the two pullbacks of \(u\) agree because the
kernel square identifies their difference with a torus translation.
The equalizer description of sections for this faithfully flat
cover gives \(q^*a=u\).  The abelian function theorem makes \(a\)
a pullback from \(S\).  Evaluation at the identity section makes that
base function one.  Thus \(u=1\).  This is precisely the route in
\srcid{even_cubic_unit_eq_one}: the mixed ratio is a character before
the restriction of \(u\) itself is made into a character.
\end{proof}

\subsection{Quotient line isomorphisms and their normalization}
Let \(N,L,M\) be invertible sheaves on \(Y,A,B\), and let
\(F:Y\to A\), \(p:Y\to B\).  Suppose the actual generating
isomorphisms are
\[
                       u:N\longrightarrow F^*L,\qquad
                       v:N\longrightarrow p^*M.
\]
For chosen frames of \(N,L,M\), write their scalar coefficients
\(u_{km}\) and \(v_{jm}\), with \(m\) the intermediate \(N\)-frame.

\begin{srlemma}[The exact quotient coefficient]
\label{lem:preunif-quotient-coefficient}
The local units \(v_{jm}u_{km}^{-1}\) define a generating isomorphism
\(w:F^*L\to p^*M\), characterized by
\[
                      v_{jm}=w_{jk}u_{km}.
 \tag{PUF4}\label{eq:preunif-quotient-coefficient}
\]
\end{srlemma}
\begin{proof}
On changing the \(N\)-frame, both \(u_{km}\) and \(v_{jm}\)
are multiplied by the same transition unit.
Their quotient is unchanged.  The chosen \(N\)-charts cover the
intersection of the pulled-back \(L\)- and \(M\)-charts, so the
quotients glue there.  On changing the \(L\)- or \(M\)-frame they
transform by exactly the transition units of
\(\operatorname{Hom}(F^*L,p^*M)\).
They therefore define \(w\).  Every quotient is a unit, so \(w\)
is generating and invertible.
Multiplication by \(u_{km}\) gives the characterization.
This expands \srcid{exists_isQuotIso} and
\srcid{IsQuotIso.isGen}.
\end{proof}

Suppose now the three lines carry the specified rigidifications,
symmetries and cubical generators, and \(u,v\) preserve those data.
If \(W(x)\) denotes the coefficient of \(w\) at a point \(x\),
cancellation of the common \(N\)-coefficients gives the identities
\[
 \begin{aligned}
 r_M&=W(1)r_L,\\
 s_M(px)W(x)&=W(x^{-1})s_L(Fx),\\
 C_M(px,py,pz)\,W(xy)W(xz)W(yz)W(1)
   &=W(xyz)W(x)W(y)W(z)\,C_L(Fx,Fy,Fz).
 \end{aligned}
 \tag{PUF5}\label{eq:preunif-quotient-cube}
\]
Here all coefficients are evaluated in the same local chart system
on the test scheme.  For the first identity divide the two
rigidification identities for \(u,v\).  For the second divide their
symmetry identities.  For the third divide their seven-point
cubical identities; the intermediate coefficients cancel
point by point.  This proves \srcid{IsQuotIso.vals} with its factor
\(W(1)\) in the stated position.

\begin{srlemma}[Uniqueness and compatibility of the normalized coefficients]
\label{lem:preunif-coefficient-unique}
For a connected family with the actual quotient
\eqref{eq:preunif-torus-extension}, two generating maps between the
same specified cubical rigidified symmetric lines, preserving all
those structures and lying over the same group map, have the same
coefficients.
The quotient coefficients in Lemma~\ref{lem:preunif-quotient-coefficient}
therefore commute with a transition between two such levels.
\end{srlemma}
\begin{proof}
The quotient of two generating maps glues to a global unit \(a\)
on the source.  To see this without choosing a global frame, divide
their coefficients on each pair of source and target charts.
Both transition factors cancel, so these local ratios agree and
cover the source.  Dividing their rigidification, symmetry and
cubical identities gives exactly \eqref{eq:preunif-even-cube}
for \(a\).  Lemma~\ref{lem:preunif-even-cubic-unit} makes \(a=1\).

For the compatibility assertion, pull the upper-level maps to the
lower level and compose with the specified source transition.
They and the given lower-level maps are over the same group map.
Their cubical comparisons are the pullbacks of the same comparisons;
uniqueness just proved identifies their coefficients.
In \eqref{eq:preunif-quotient-coefficient}, cancel the common
coefficient of the intermediate line to identify the two
quotient coefficients.  This is \srcid{gw_eq_of_isCPMap}
followed by \srcid{gw_compat}.
\end{proof}

\subsection{The finite translation-kernel argument over a local ring}
\begin{srlemma}[Finite augmented algebra with trivial closed fibre]
\label{lem:preunif-finite-augmentation}
Let \((A,\mathfrak m)\) be local, let \(C\) be a finite \(A\)-algebra,
and let \(\epsilon:C\to A\) be a retraction of its structural map
\(\phi:A\to C\).  Suppose
\[
                    c-\phi(\epsilon c)\in\mathfrak m C
                    \qquad(c\in C).
 \tag{PUF6}\label{eq:preunif-augmentation-reduction}
\]
Then \(\phi\epsilon=1_C\); in particular \(C\simeq A\).
\end{srlemma}
\begin{proof}
The finite \(A\)-module \(C/\phi(A)\) is generated modulo
\(\mathfrak m\) by zero, by
\eqref{eq:preunif-augmentation-reduction}.
Nakayama's lemma makes it zero.  Thus \(c=\phi(a)\) for some \(a\).
Applying the retraction gives \(\epsilon c=a\), and hence
\(\phi(\epsilon c)=c\).
No flatness of \(C\) is assumed.  This is
\srcid{ring_retraction}.
\end{proof}

\begin{srproposition}[The exact closed-fibre translation-kernel deduction]
\label{prop:preunif-ktriv-local}
Let \(B/A\) be an abelian scheme over a local Noetherian ring,
and let \(L\) be the specified relatively ample line.
Retain the actual closed translation-kernel representation:
a finite closed subgroup \(j:H\hookrightarrow B\) represents
\(\operatorname{InK}(B,L)\) on every test scheme.
If the pullback of \(L\) to the residue-field fibre has trivial
translation kernel on every test scheme, then
\(\operatorname{KTriv}(B,L)\) holds over \(A\).
\end{srproposition}
\begin{proof}
The identity translation gives a section \(e:\operatorname{Spec}A
\to H\).  Because \(H\) is finite, it is \(\operatorname{Spec}C\)
for a finite \(A\)-algebra.  The section gives the augmentation
\(\epsilon:C\to A\).

Any point of \(H\) over a residue-field test scheme maps, through
the Cartesian fibre diagram and \(j\), to a translation preserving
the special-fibre line.  The hypothesis makes it the identity.
Since \(j\) is a closed immersion and hence a monomorphism,
the point of \(H\) itself equals the identity section.
This assertion is for all test schemes: in particular it applies
to the universal point of
\(\operatorname{Spec}(C/\mathfrak mC)\).
On coordinate rings it is exactly
\eqref{eq:preunif-augmentation-reduction}.
Lemma~\ref{lem:preunif-finite-augmentation} makes
\(H=\operatorname{Spec}A\) with its identity section.
Every \(\operatorname{InK}\)-point factors through \(H\) and is
therefore the identity, as required.
Checking only geometric points of the closed fibre would not
justify this argument, since it would discard nilpotents.

In the source, the representation is
\srcid{exists_closed_twist_kernel}, its finiteness is
\srcid{AmpleStabFinite.isFinite_of_relAmple}, and
\srcid{inK_iff_trivial_twist} supplies its stated equality with
\(\operatorname{InK}\).
This proposition prints their subsequent local deduction,
\srcid{kTriv_of_closed_fibre}; it retains those representation
and finiteness producers rather than replacing them by a claim
about geometric point counts.
\end{proof}

\subsection{Compatible weight points and their exact algebraisation}
Let \(R\) now be a Noetherian local domain, complete and separated
for the proper ideal \(I\), and put \(K=\operatorname{Frac}(R)\).  Put \(R_e=R/I^{e+1}\).
Let \(G_e\) be the actual pullback of the semiabelian family \(G/R\),
and suppose the formal quotient tower has quotients
\(q_e:G_e\to B_e\) by the same character lattice \(X\), cubical lines
\(M_e\), transitions \(\beta_e:B_e\to B_{e+1}\), and their specified
Cartesian and torus-compatibility squares.
Retain its actual algebraisation \(a\): an abelian scheme \(B/R\),
a relatively ample cubical rigidified symmetric line \(M\), and
compatible Cartesian maps \(b_e:B_e\to B\).
Assume \(\operatorname{KTriv}(B,M)\), for example from
Proposition~\ref{prop:preunif-ktriv-local} with its stated inputs.

The source's representation condition
\(\operatorname{Rep}(e,\chi,x_e)\) means that \(x_e\) lies over
\(\operatorname{Spec}R_e\), and the actual translation-ratio line
\[
                  K_{x_e}(t)=M_{x_e+t}^{-1}\otimes M_t,
\]
pulled through the specified \(b_e q_e\), has a generating section
with torus weight \(\chi\). The two line exponents are
\([-1,+1]\), as in the source definitions
\srcid{Rep} and \srcid{TowerR.Alg.RepWt}.  Its uniqueness uses
\(\operatorname{KTriv}\), and its existence is the precise
Picard-zero degeneration input
\(\operatorname{PicardZeroDegenerationCondition}\).
In the following theorem that existence input is stated explicitly.

\begin{srlemma}[Algebraising a compatible family of points]
\label{lem:preunif-adic-points}
Let \(Y/R\) be a scheme and let
\(x_e:\operatorname{Spec}R_e\to Y\) be compatible \(R\)-maps.
There is a unique \(R\)-map \(x:\operatorname{Spec}R\to Y\)
whose reductions are all the \(x_e\).
\end{srlemma}
\begin{proof}
All closed points of these local schemes have the same image in
\(Y\).  Choose an affine open \(U\) containing that image.
The preimage of \(U\) under any \(x_e\) is an open set containing
the closed point of \(\operatorname{Spec}R_e\), so is the whole
spectrum.  The compatible coordinate maps therefore give
\[
                \Gamma(U,\mathcal O)\longrightarrow
                  \varprojlim_e R_e=R.
\]
This defines \(x\), and compatibility with the structural map
makes it an \(R\)-map.

For uniqueness, two such maps have the same closed-point image.
Both factor through the chosen affine \(U\), for the same
local-spectrum reason.  Their coordinate maps agree modulo every
\(I^{e+1}\), and separatedness of the \(I\)-adic topology makes
them equal.  This is the actual local argument behind
\srcid{existsUnique_over_of_isAdicFamily}.
\end{proof}

\begin{srproposition}[The additive character-point map]
\label{prop:preunif-additive-weights}
Assume that, for each \(e\) and \(\chi\in X\), there exists an
\(x_{e,\chi}\) satisfying the actual representation condition above.
Then these points give one additive map
\[
                         c:X\longrightarrow B(R)
\]
with the source's \(\operatorname{RepWt}(c_\chi,\chi,e)\)
for every \(e\).
\end{srproposition}
\begin{proof}
Pulling a representing point and its weighted generating section
through a tower transition preserves its character: this is the
actual torus-compatibility square.  The two resulting points at
level \(e\) have the same character representation.
Their difference is in the translation kernel of \(M\);
\(\operatorname{KTriv}\) makes them equal.
Thus the points \(x_{e,\chi}\), after their maps into \(B\), are
compatible in \(e\).  Lemma~\ref{lem:preunif-adic-points} supplies
\(c_\chi\).

For additivity use the specified cube.  For points \(a,b,t\), it
identifies the square line
\[
          S_{a,b}(t)
               =M_{a+b+t}M_tM_{a+t}^{-1}M_{b+t}^{-1}
\]
with the constant line \(M_{a+b}M_0M_a^{-1}M_b^{-1}\) from the
base.  Over the local ring this constant invertible module is
free.  Choose a generating section, and pull it back to the
actual torus quotient.  Its torus weight is zero, since the
quotient map is invariant under the kernel action.
The tensor identity, with every exponent shown, is
\[
                    K_a\otimes K_b\otimes S_{a,b}^{-1}
                                \simeq K_{a+b}.
\]
Indeed its coefficient line is
\[
 (M_{a+t}^{-1}M_t)(M_{b+t}^{-1}M_t)
       (M_{a+b+t}^{-1}M_t^{-1}M_{a+t}M_{b+t})
                    =M_{a+b+t}^{-1}M_t .
\]
Tensor the two weight generators with the inverse square generator.
Their weights are \(\chi,\chi'\) and zero, so the resulting
generator has weight \(\chi+\chi'\). This is exactly
\srcid{triv_square_sections}, \srcid{Square.triv_inv} and
\srcid{exists_genWt_add}; no character is dualized.
Therefore \(c_\chi+c_{\chi'}\) and \(c_{\chi+\chi'}\) satisfy the
same representation condition at every level.
Uniqueness at every level, followed by the uniqueness part of
Lemma~\ref{lem:preunif-adic-points}, identifies the two sections.
This is \srcid{exists_repWt}, \srcid{repWt_add} and
\srcid{repWt_unique}, giving \srcid{exists_weightPoints}.
\end{proof}

\subsection{Assembly of the actual preliminary datum}
Retain the preceding tower and algebraisation, its character map
\(c\), the specified line \(P\) on \(G\), and an actual generic
polarized abelian variety \(A/K\) with the fixed group isomorphism
\[
                    G_K\simeq A
\]
and its specified generating comparison of generic lines.
Let the source torus-coordinate producer be applied to the actual
\(q_e\), the actual \(b_e\), and the weight generators just obtained.
It gives coordinates \(g_e\) with the same torus embedding and the
same map
\[
                        G_e\xrightarrow{q_e}B_e
                                  \xrightarrow{b_e}B .
\]
Its coordinate-compatibility assertion
\srcid{torusCoordinates_compat_of_torus_compat} identifies the
generators at consecutive levels, once the quotient maps and the
torus embeddings agree.

\begin{srproposition}[The source's \(\mathrm{PreU}\) assembly]
\label{prop:preunif-assembly}
Under these same tower, algebraisation, kernel, weight and
coordinate inputs, there are an actual preliminary degeneration
datum \(D\) and \(\operatorname{PreU}(I,D,A,G)\), whose line on
\(G\) and generic isomorphism are the originally specified ones.
\end{srproposition}
\begin{proof}
Take \(D.P\) to be \(B\) with line \(M\), its actual
rigidification, symmetry and cube.  The abelian relative dimension
is the constant dimension of \(B/R\), since the local base is
connected.  The relative ampleness and \(\operatorname{KTriv}\)
are the already specified properties of this algebraisation.
Take \(D.X=X\) and its point map to be \(c\).
For each \(\chi\), choose a generic generator of the inverse
one-dimensional fibre \(M(c_\chi)^{-1}\), over the field \(K\).
The exponent is \(-1\), exactly as in the source's
\srcid{DVR.exists_gen_of_local} application.
This is an initial choice of \(\psi_\chi\), not yet the
nonzero-theta-normalized choice needed for the full degeneration
datum.

At level \(e\), let \(P_e\) be the actual pullback of \(P\).
It has one generating comparison with \(P|_{G_e}\), and its
descended line has the actual generating comparison with
\(M\) along \(b_eq_e\).
Apply Lemma~\ref{lem:preunif-quotient-coefficient} to obtain the
generating comparison
\[
                   P|_{G_e}\longrightarrow (b_eq_e)^*M.
\]
Its rigidification, symmetry and cubical values are
\eqref{eq:preunif-quotient-cube}.
The quotient maps commute with the tower transitions, by the
tower equation and \(b_{e+1}\beta_e=b_e\).
The torus embeddings commute by the original torus-compatibility
condition.  Thus the \(g_e\) coordinate generators commute, by
the coordinate producer just specified.
Finally Lemma~\ref{lem:preunif-coefficient-unique} makes the
line-comparison coefficients commute too.
The needed connectedness at each level holds because \(R_e\)
is local and the semiabelian family is smooth with geometrically
connected fibres.  We have supplied exactly the maps, generators,
value identities and level compatibilities in
\(\operatorname{PreU}\).
The prescribed line on \(G\) and the prescribed generic
isomorphism were kept throughout, so their equalities are
literal.  This is the construction \srcid{exists_preU}.
\end{proof}

\paragraph{The tower and algebraisation frontier, with its original order.}
The source obtains the actual \(\operatorname{TowerR}\) as follows.
Start with the given split extension over \(R/I\), with its free
finite character lattice \(X\).
Inductively lift its torus quotient through
\(R/I^{e+2}\to R/I^{e+1}\), whose kernel is nilpotent.
At each level, trivialise the cubical biextension first over the
normal quotient \(R/I\), lift that trivialisation across the
nilpotent ideal, and use the specified
\(\operatorname{HypW}\) relation to obtain the even cubical
trivialisation.  Descent supplies the quotient cubical line.
Its transition is the descent of the actual compatible torus
quotient maps.  These are the successive producers
\srcid{exists_tq_succ}, \srcid{nonempty_biTriv},
\srcid{exists_even}, and \srcid{exists_isCPPullback} in
\srcid{NonemptyFormalUniformisationTower}.
They are not replaced by the assertion that compatible
quotients of schemes automatically exist.

Next the actual closed-fibre polarization producer proves
relative ampleness and \(\operatorname{KTriv}\) on each
field fibre of the first quotient.  The relative-ampleness
producer over the complete local principal quotient gives
relative ampleness at level zero.
The project's \(\operatorname{Ega}\) theorem algebraises this
ample formal tower to a projective smooth scheme with its line.
The group multiplication, inverse and identity are algebraised
from their compatible level maps; equality of the group laws
is checked on all levels using formal faithfulness.
The original producer
\srcid{IsAbelianSchemeOfAbelianThickenings} additionally proves
geometric connectedness.  Then
\srcid{exists_cp_alg} algebraises the actual cubical, symmetric
and rigidified line structures, rather than just its isomorphism
class.
This is \srcid{NonemptyAlgebraisationFormalTowerAmple}.
The local kernel argument above is applied first to the
initial quotient and then to this actual algebraisation.
The Picard-zero degeneration input then supplies the
representing weight points used in
Proposition~\ref{prop:preunif-additive-weights}.

These lower existence and algebraisation producers are the
remaining construction inputs of this section.  The full N3
argument must also replace the initially chosen generic
\(\psi_\chi\) by the generators obtained from the same nonzero
theta section, compare with the chosen complete-discrete-
valuation-ring uniformisation, descend its expansion
bijections and then construct the actual linear PEL/torsion
identifications.  Those steps are not consequences of the
coefficient and local-kernel arguments proved here.

\paragraph{Frozen source correspondence.}
All the following paths are in
\srcid{Development/AbelianVarietyTheory/Proofs/} at commit
\srcid{7cbf7a0c32fca1badf05931879a37c3f8a2b03b4}:
\srcid{UniformizationProofs/EvenCubicUnitEqOne.lean},
\srcid{DegenerationProofs/ExistsPreUOfFibres.lean},
\srcid{DegenerationProofs/KTrivOfClosedFibreLocalRing.lean},
\srcid{DegenerationProofs/ExistsAdditiveWeightPointsOfTorusTower.lean},
\srcid{DegenerationProofs/NonemptyFormalUniformisationTower.lean},
and
\srcid{DegenerationProofs/NonemptyAlgebraisationFormalTowerAmple.lean}.
The last paragraph records the exact remaining source producers;
it is not a claim that all of their proofs have already been
printed in this section.
